Fix a sufficiently large index and choose, using the additional premise in the proposed statement, a member of \(\mathcal M_{n,d,L}\) for which the latent types are independent uniform draws.  Retain its kernel \(W\), its spectral data, and its graph law throughout the proof.  Put
\[
 w(u)=\int_0^1W(u,v)\,dv,
 \qquad
 p(u)=\rho_n w(u)=\frac dn w(u),
 \qquad
 \mu(u)=(n-1)p(u).
\tag{1}
\]
By \({\rm ass:kernel\text{-}bound}\) and \({\rm ass:degree\text{-}regularity}\),
\[
 \kappa^{-1}\leq w(u)\leq\kappa
\tag{2}
\]
for Lebesgue-almost every \(u\).  The lower inequality in (2) is the stated row-regularity condition, while the upper inequality follows by integrating \(W(u,v)\leq\kappa\) over a set of measure one.  Because \(U_i\) is uniform, (2) holds at every sampled \(U_i\) with probability one.  The upper bound \(p(u)\leq1\) follows from \({\rm ass:edge\text{-}probability}\).

Conditional on \(U_i=u\), the pairs consisting of \(U_j\) and the edge randomization for \(A_{ij}\), for \(j\ne i\), are independent by \({\rm ass:unit\text{-}sampling}\) and \({\rm ass:graphon\text{-}draw}\).  Integrating one such pair over the uniform \(U_j\) gives
\[
 \Pr(A_{ij}=1\mid U_i=u)
 =\rho_n\int_0^1W(u,v)\,dv=p(u).
\tag{3}
\]
Thus the conditional degree law is
\[
 N_i\mid U_i=u\sim{\rm Binomial}(n-1,p(u)).
\tag{4}
\]
For \(n\geq2\), equations (1)--(2) imply the uniform mean bounds
\[
 \frac{d}{2\kappa}\leq\mu(u)\leq\kappa d.
\tag{5}
\]
This is the only graphon-specific input to the testing constructions below.

\textit{Step 1: class-valid response pairs.}
Let \(s_0=B/4>0\) and let \(\xi_i\), independently across units and independently of \((U,A,Z)\), have density
\[
 p_{s_0}(t)=s_0^{-1}\cos^2\!\left(\frac{\pi t}{2s_0}\right)
 \mathbf 1\{|t|\leq s_0\}.
\tag{6}
\]
The substitution \(y=\pi t/(2s_0)\) gives
\[
 \int_{\mathbb R}p_{s_0}(t)\,dt
 =s_0^{-1}\int_{-s_0}^{s_0}
 \cos^2\!\left(\frac{\pi t}{2s_0}\right)dt=1,
\tag{7}
\]
so (6) is a probability density, and \(|\xi_i|\leq B/4\) almost surely.  Given a deterministic \(g\in C^3[0,1]\), define two response laws by
\[
 f_i^\pm(a,x)=\xi_i\pm g(x),
 \qquad a\in\{0,1\},\quad x\in[0,1].
\tag{8}
\]
Suppose that
\[
 \|g\|_\infty\leq\frac{3B}{4},
 \qquad
 \|g'\|_\infty\leq B,
 \qquad
 \|g''\|_\infty\leq B,
 \qquad
 \|g'''\|_\infty\leq L.
\tag{9}
\]
Then (6), (8), and (9) verify \({\rm ass:response\text{-}envelope}\) and \({\rm ass:third\text{-}derivative\text{-}bound}\).  Equation (8), inserted into the exposure formula in \({\rm ass:anonymous\text{-}interference}\), defines the corresponding potential outcomes.  The pairs \((U_i,f_i^\pm)\) are independent and identically distributed, so \({\rm ass:unit\text{-}sampling}\) is preserved.  The distribution of \(U_i\) and every graphon eigenfunction are unchanged, hence the Bernstein condition in \({\rm ass:eigenfunction\text{-}bernstein}\) is inherited from the selected member.  All kernel, rank, spectral-strength, graph, and design conditions are unchanged.  Consequently, whenever (9) holds, both laws generated by (8) belong to \(\mathcal M_{n,d,L}\).

The derivative of the random intercept in (8) is zero and the same \(g\) is used for both values of \(a\).  Therefore \({\rm def:target}\) gives the deterministic target values
\[
 \theta_n^\pm=\pm g'(1/2),
 \qquad
 \Delta_g:=|\theta_n^+-\theta_n^-|=2|g'(1/2)|.
\tag{10}
\]

\textit{Step 2: a common distance bound.}
Let \(q_{s_0}=\sqrt{p_{s_0}}\), extended by zero outside \([-s_0,s_0]\).  Since the cosine in (6) is nonnegative on its support and vanishes at both endpoints, this extension is absolutely continuous.  Direct integration gives
\[
 \|q_{s_0}'\|_2^2
 =\frac1{s_0}\left(\frac{\pi}{2s_0}\right)^2
 \int_{-s_0}^{s_0}\sin^2\!\left(\frac{\pi t}{2s_0}\right)dt
 =\frac{\pi^2}{4s_0^2}.
\tag{11}
\]
For translations \(\eta_1,\eta_2\), the fundamental theorem of calculus and Minkowski's integral inequality yield
\[
 \begin{aligned}
 H^2\{p_{s_0}(\mathord\cdot-\eta_1),p_{s_0}(\mathord\cdot-\eta_2)\}
 &=\|q_{s_0}(\mathord\cdot-\eta_1)-q_{s_0}(\mathord\cdot-\eta_2)\|_2^2\\
 &\leq(\eta_1-\eta_2)^2\|q_{s_0}'\|_2^2
 =\frac{\pi^2(\eta_1-\eta_2)^2}{4s_0^2}.
 \end{aligned}
\tag{12}
\]
Here \(H^2\) is the integral of the squared difference of the square-root densities.  Define the observed exposure argument
\[
 X_i=\begin{cases}
 1/2+V_i,&N_i>0,\\
 0,&N_i=0.
 \end{cases}
\tag{13}
\]
Conditional on \((A,Z)\), the outcomes under the two signs in (8) are independent translates of (6), with translations separated by \(2g(X_i)\).  For product measures, square-root affinities multiply.  Since \(1-\prod_i(1-x_i)\leq\sum_i x_i\) for \(0\leq x_i\leq1\), (12) and averaging over the common law of \((A,Z)\) give
\[
 H^2(Q_g^+,Q_g^-)
 \leq\frac{\pi^2}{s_0^2}\mathbb E\sum_{i=1}^ng(X_i)^2,
\tag{14}
\]
where \(Q_g^+\) and \(Q_g^-\) are the complete observed-data laws.  Cauchy--Schwarz applied to the two square-root densities also gives
\[
 \operatorname{TV}(Q_g^+,Q_g^-)
 \leq H(Q_g^+,Q_g^-).
\tag{15}
\]

\textit{Step 3: uniform degree and local-mass controls.}
For \(N\sim{\rm Binomial}(n-1,p)\), expansion of the binomial polynomial and integration of \(t^N\) on \([0,1]\) give the exact identity
\[
 \mathbb E\frac1{N+1}
 =\int_0^1(1-p+pt)^{n-1}dt
 =\frac{1-(1-p)^n}{np}
 \leq\frac1{np}.
\tag{16}
\]
Because \(N^{-1}\mathbf1\{N>0\}\leq2/(N+1)\), equations (2)--(4) and (16) imply
\[
 \mathbb E\{N_i^{-1}\mathbf1(N_i>0)\mid U_i=u\}
 \leq\frac{2}{np(u)}\leq\frac{2\kappa}{d}.
\tag{17}
\]
Conditional on \((A,N_i=N>0)\), \({\rm ass:bernoulli\text{-}design}\) gives
\[
 M_i\sim{\rm Binomial}(N,1/2),
 \qquad
 \mathbb E_Z(V_i^2\mid A,N_i=N)=\frac1{4N}.
\tag{18}
\]
Integrating (17)--(18) yields
\[
 \mathbb E\sum_{i=1}^nV_i^2\mathbf1\{N_i>0\}
 \leq\frac{\kappa n}{2d}.
\tag{19}
\]
The lower mean bound in (5) also gives
\[
 \Pr(N_i=0\mid U_i=u)=(1-p(u))^{n-1}
 \leq e^{-\mu(u)}\leq e^{-d/(2\kappa)}.
\tag{20}
\]

We next bound the probability of a localized exposure.  For \(N=2q\), put \(a_q={2q\choose q}4^{-q}\).  Then
\[
 \frac{a_{q+1}}{a_q}=\frac{2q+1}{2q+2},
 \qquad
 \left(\frac{2q+1}{2q+2}\right)^2\leq\frac{q+1}{q+2}.
\tag{21}
\]
Starting from \(a_0=1\), induction in (21) gives \(a_q\leq(q+1)^{-1/2}\).  The maximal atom for an odd number of trials is no larger than the adjacent even maximal atom.  Hence there is a numerical constant \(C_0\) such that every atom of a \({\rm Binomial}(N,1/2)\) law is at most \(C_0N^{-1/2}\).  The interval \(|M/N-1/2|\leq h\) contains at most \(2Nh+2\) integer values, so
\[
 \Pr\{|M/N-1/2|\leq h\mid N\}
 \leq C_1(\sqrt N\,h+N^{-1/2}).
\tag{22}
\]

For the mixture (4), conditional Jensen gives
\[
 \mathbb E(\sqrt{N_i}\mid U_i=u)\leq\sqrt{\mu(u)}\leq\sqrt{\kappa d}.
\tag{23}
\]
To control the inverse square root, exponential Markov with \(2^{-N_i}\) gives
\[
 \begin{aligned}
 \Pr\{N_i<\mu(u)/2\mid U_i=u\}
 &\leq 2^{\mu(u)/2}\mathbb E(2^{-N_i}\mid U_i=u)\\
 &=2^{\mu(u)/2}(1-p(u)/2)^{n-1}\\
 &\leq\exp\{-c_0\mu(u)\},
 \end{aligned}
 \qquad c_0=\frac{1-\log2}{2}>0.
\tag{24}
\]
Splitting at \(\mu(u)/2\), then using (5), shows, for a constant \(C_2\) depending only on \(\kappa\),
\[
 \begin{aligned}
 \mathbb E\{N_i^{-1/2}\mathbf1(N_i>0)\mid U_i=u\}
 &\leq\sqrt{2/\mu(u)}+e^{-c_0\mu(u)}\\
 &\leq C_2d^{-1/2}
 \end{aligned}
\tag{25}
\]
for all sufficiently large \(n\).  Here \(d\to\infty\) follows from \({\rm ass:log\text{-}degree}\).  Combining (18), (22)--(25), and \(h\geq d^{-1}\) from \({\rm def:bandwidth\text{-}set}\), we obtain, uniformly over \(h\in\mathcal H_{n,d}\),
\[
 \Pr(|V_i|\leq h,N_i>0)
 \leq C_3\{\sqrt d\,h+d^{-1/2}\}
 \leq2C_3\sqrt d\,h.
\tag{26}
\]

Finally, choose \(K>e\kappa\) so large that
\[
 a_K:=K\log\!\left(\frac{K}{e\kappa}\right)
 \quad\hbox{satisfies}\quad a_KC_{\log}>2,
\tag{27}
\]
and put \(D_n=\lceil Kd\rceil\).  For a binomial variable with mean \(\mu\leq\kappa d\), exponential Markov optimized at \(t=\log(D_n/\mu)>0\) gives
\[
 \Pr(N_i>D_n\mid U_i=u)
 \leq\left(\frac{e\mu(u)}{D_n}\right)^{D_n}
 \leq e^{-a_Kd}.
\tag{28}
\]
The union bound, (27), and \({\rm ass:log\text{-}degree}\) now imply
\[
 \Pr\left(\max_{1\leq i\leq n}N_i>D_n\right)
 \leq ne^{-a_Kd}
 \leq n^{1-a_KC_{\log}}=o(1).
\tag{29}
\]
No independence among the degrees is used in (29).

\textit{Step 4: the three least-favourable response pairs.}
For the graph term, let
\[
 g_G(x)=c_GB\sqrt{d/n}\,(x-1/2).
\tag{30}
\]
The sparse-limit condition gives \(d/n\to0\), so (9) holds for every sufficiently large \(n\) whenever \(0<c_G\leq1\).  Equations (13), (19), (20), and (30) give
\[
 \begin{aligned}
 \mathbb E\sum_{i=1}^ng_G(X_i)^2
 &=c_G^2B^2\frac dn
 \left\{\mathbb E\sum_iV_i^2\mathbf1(N_i>0)
       +\frac n4\Pr(N_i=0)\right\}\\
 &\leq c_G^2B^2\left\{\frac\kappa2+\frac d4e^{-d/(2\kappa)}\right\}
 \leq C_4c_G^2B^2.
 \end{aligned}
\tag{31}
\]
Let
\[
 \eta=\frac{1-2\alpha}{2}>0.
\tag{32}
\]
Since \(s_0=B/4\), equations (14)--(15) and (31) permit a fixed choice of \(c_G>0\), depending only on the fixed class constants and \(\alpha\), for which
\[
 \operatorname{TV}(Q_G^+,Q_G^-)\leq\eta.
\tag{33}
\]
By (10) and (30), the target separation is
\[
 \Delta_G=2c_GB\sqrt{d/n}.
\tag{34}
\]

For the interior term, suppose \(L>0\) and define
\[
 \phi(t)=t(1-4t^2)^4\mathbf1\{|t|<1/2\}.
\tag{35}
\]
The polynomial in (35) has a zero of order four at each endpoint, so its extension by zero is in \(C^3(\mathbb R)\), and \(\phi'(0)=1\).  Write \(C_j=\|\phi^{(j)}\|_\infty\) for \(0\leq j\leq3\).  For \(h\in\mathcal H_{n,d}\), set
\[
 g_h(x)=c_ILh^3\phi\!\left(\frac{x-1/2}{h}\right).
\tag{36}
\]
Then
\[
 \|g_h^{(j)}\|_\infty=c_ILh^{3-j}C_j,
 \qquad 0\leq j\leq3.
\tag{37}
\]
Choose \(c_IC_3\leq1\).  Since \(h\leq d^{-1/2}\to0\) and \(B,L\) are fixed, all the other inequalities in (9) follow from (37) for sufficiently large \(n\).  The bump vanishes at the isolate exposure zero and is supported only when \(|V_i|\leq h\).  Therefore (14), (26), and (36) give
\[
 H^2(Q_h^+,Q_h^-)
 \leq C_5c_I^2n\sqrt d\,L^2h^7.
\tag{38}
\]
Define the proof-intermediate bandwidth
\[
 h_0=\left(\frac{L^{-2}}{n\sqrt d}\right)^{1/7}.
\tag{39}
\]
When \(d^{-1}<h_0<d^{-1/2}\), substitution into (38), followed if necessary by a fixed decrease of \(c_I\), gives
\[
 \operatorname{TV}(Q_{h_0}^+,Q_{h_0}^-)\leq\eta.
\tag{40}
\]
Equations (10), (36), and (39) yield
\[
 \Delta_I=2c_ILh_0^2
 =2c_IL^{3/7}(n^2d)^{-1/7}.
\tag{41}
\]

For the lattice term, retain \(D_n\) from (27)--(29), put \(w_n=(4D_n)^{-1}\), and define
\[
 g_D(x)=c_D Lw_n^3\phi\!\left(\frac{x-1/2}{w_n}\right),
\tag{42}
\]
where \(c_DC_3\leq1\).  The analogue of (37) and \(D_n\to\infty\) verify (9) for all sufficiently large \(n\).  If \(1\leq N_i\leq D_n\), every noncentral point of the exposure lattice satisfies
\[
 |V_i|\geq
 \begin{cases}
 N_i^{-1},&N_i\text{ is even},\\
 (2N_i)^{-1},&N_i\text{ is odd},
 \end{cases}
 \geq(2D_n)^{-1}=2w_n.
\tag{43}
\]
Such a point lies outside the support of (42), while oddness gives \(g_D(1/2)=0\) at the central point and the support condition gives \(g_D(0)=0\) for isolates.  Hence on
\[
 E_n=\left\{\max_{1\leq i\leq n}N_i\leq D_n\right\}
\tag{44}
\]
the observed outcomes under both signs equal \(\xi_i\).  The restrictions of the two observed-data laws to \(E_n\) are identical, so (29) gives
\[
 \operatorname{TV}(Q_D^+,Q_D^-)
 \leq\Pr(E_n^c)=o(1)\leq\eta
\tag{45}
\]
for all sufficiently large \(n\).  On the other hand, (10) and (42) give
\[
 \Delta_D=2c_DLw_n^2\asymp Ld^{-2},
\tag{46}
\]
where both comparison constants depend only on \(K\), hence only on fixed class constants.

\textit{Step 5: reduction to minimax risk and honest expected length.}
Consider any one of the preceding pairs and abbreviate its target values by \(\theta^+,\theta^-\), its data laws by \(Q^+,Q^-\), and their separation by \(\Delta\).  For an arbitrary estimator \(T\), let \(A_T=\{|T-\theta^+|\geq\Delta/2\}\).  On \(A_T^c\), the triangle inequality gives \(|T-\theta^-|\geq\Delta/2\).  Therefore
\[
 \begin{aligned}
 \mathbb E_{Q^+}|T-\theta^+|+\mathbb E_{Q^-}|T-\theta^-|
 &\geq\frac\Delta2\{Q^+(A_T)+Q^-(A_T^c)\}\\
 &\geq\frac\Delta2\{1-\operatorname{TV}(Q^+,Q^-)\}.
 \end{aligned}
\tag{47}
\]
Thus the larger of the two risks is at least \(\Delta\{1-\operatorname{TV}(Q^+,Q^-)\}/4\).

Let \(\mathcal C\) be any confidence interval admissible in \({\rm def:minimax\text{-}criteria}\).  Honesty under both members of a pair, total variation, and the union bound imply
\[
 \begin{aligned}
 Q^+\{\theta^+,\theta^-\in\mathcal C\}
 &\geq Q^+\{\theta^+\in\mathcal C\}
       +Q^+\{\theta^-\in\mathcal C\}-1\\
 &\geq(1-\alpha)+(1-\alpha-\operatorname{TV}(Q^+,Q^-))-1\\
 &=1-2\alpha-\operatorname{TV}(Q^+,Q^-).
 \end{aligned}
\tag{48}
\]
On this event the interval length is at least \(\Delta\).  Equations (32)--(33), (40), and (45) therefore give
\[
 \mathbb E_{Q^+}\operatorname{length}(\mathcal C)
 \geq\Delta\frac{1-2\alpha}{2}
\tag{49}
\]
for each applicable pair.  Taking the infima in \({\rm def:minimax\text{-}criteria}\), (34), (41), and (46) show that each of \(R^*_{n,d,L}\) and \(\Lambda^*_{n,d,L}(\alpha)\) is bounded below, up to fixed positive constants, by \(\sqrt{d/n}\), by \(L^{3/7}(n^2d)^{-1/7}\) when (39) is interior, and by \(Ld^{-2}\), respectively.

\textit{Step 6: comparison with the supported envelope.}
Write \(a_n=(n\sqrt d)^{-1/2}\).  If \(L>0\), equation (39) gives
\[
 Lh_0^2=a_nh_0^{-3/2}=L^{3/7}(n^2d)^{-1/7}.
\tag{50}
\]
If \(h_0\leq d^{-1}\), then \(L^2\geq d^{13/2}/n\), so
\[
 a_nd^{3/2}=d^{5/4}n^{-1/2}\leq Ld^{-2}.
\tag{51}
\]
Evaluating the infimum in \({\rm def:frontier\text{-}envelope}\) at \(h=d^{-1}\) gives
\[
 \inf_{h\in\mathcal H_{n,d}}\{Lh^2+a_nh^{-3/2}\}
 \leq2Ld^{-2}.
\tag{52}
\]
If \(d^{-1}<h_0<d^{-1/2}\), evaluation at \(h=h_0\) and (50) give
\[
 \inf_{h\in\mathcal H_{n,d}}\{Lh^2+a_nh^{-3/2}\}
 \leq2L^{3/7}(n^2d)^{-1/7}.
\tag{53}
\]
If \(h_0\geq d^{-1/2}\), then \(L^2\leq d^3/n\), and evaluation at \(h=d^{-1/2}\) gives
\[
 Ld^{-1}\leq\sqrt{d/n},
 \qquad
 \inf_{h\in\mathcal H_{n,d}}\{Lh^2+a_nh^{-3/2}\}
 \leq2\sqrt{d/n}.
\tag{54}
\]
When \(L=0\), the infimum is attained at \(h=d^{-1/2}\) and equals \(\sqrt{d/n}\).  In the regime of (52), combine the graph and lattice pairs; in the regime of (53), combine the graph and interior pairs; and in the regime of (54), or when \(L=0\), use the graph pair.  If a nonnegative minimax criterion is at least \(c_1x\) and at least \(c_2y\), it is at least \(\min(c_1,c_2)(x+y)/2\).  Applying this observation to (34), (41), (46), and (52)--(54) proves, for a constant \(c>0\) depending only on the fixed class constants and \(\alpha\),
\[
 R^*_{n,d,L}\geq c\left[\sqrt{d/n}+
 \inf_{h\in\mathcal H_{n,d}}\{Lh^2+(n\sqrt d\,h^3)^{-1/2}\}\right],
\tag{55}
\]
and the same inequality with \(R^*_{n,d,L}\) replaced by \(\Lambda^*_{n,d,L}(\alpha)\).  The bracket in (55) is \(\mathfrak r^{\mathrm{sup}}_{n,d,L}\) by \({\rm def:frontier\text{-}envelope}\).  The witness laws all belong to the full class, so lower bounds obtained on this uniform-latent submodel apply to the full-class minimax criteria.  This proves the proposed narrowed statement.